\documentclass[11pt,a4paper]{article}

% ---------- Packages ----------
\usepackage[T1]{fontenc}
\usepackage[utf8]{inputenc}
\usepackage{lmodern}
\usepackage[margin=2.5cm]{geometry}
\usepackage{parskip}
\usepackage{amsmath}
\usepackage{xcolor}
\usepackage{listings}
\usepackage{enumitem}
\usepackage{booktabs}
\usepackage{array}
\usepackage[hidelinks]{hyperref}

% ---------- Code listing style ----------
\lstdefinestyle{cstyle}{
  language=C,
  basicstyle=\ttfamily\small,
  keywordstyle=\color{blue!70!black}\bfseries,
  commentstyle=\color{green!50!black}\itshape,
  stringstyle=\color{red!60!black},
  numbers=left,
  numberstyle=\tiny\color{gray},
  frame=single,
  rulecolor=\color{gray!50},
  backgroundcolor=\color{gray!5},
  breaklines=true,
  tabsize=4,
  showstringspaces=false,
  xleftmargin=1.5em,
  framexleftmargin=1.5em
}
\lstset{style=cstyle}

\newcommand{\code}[1]{\texttt{#1}}

\title{\textbf{Ludo Game Report}}
\author{}
\date{}

\begin{document}

\maketitle

% =====================================================================
\section{Data Structures Used to Represent Players and Pieces}

\subsection{Piece Structure}

Each piece is represented by a structure containing three variables.

\begin{lstlisting}
typedef struct {
    short location;
    short index;
    short captureCounter;
} piece;
\end{lstlisting}

\begin{description}[style=nextline, leftmargin=1.5em]
  \item[\code{location}]
    This attribute is central to the piece's state within the game. It encodes
    the entire position logic of the piece, which can take the following values:
    \begin{itemize}[nosep]
      \item $-2$: the piece is in the base.
      \item $-1$: the piece has reached home.
      \item $0$ to $51$: positions on the standard path.
      \item $-10$ to $-5$: positions in the home straight.
    \end{itemize}

  \item[\code{index}]
    Identifies the piece among the four pieces belonging to the player, allowing
    easy access to and manipulation of a specific piece when needed.

  \item[\code{captureCounter}]
    Keeps track of how many captures the piece has made. This information is
    required when determining whether the piece can enter the home straight.
\end{description}

\subsection{Player Structure}

Each player is represented by a structure containing seven variables.

\begin{lstlisting}
typedef struct {
    char  color[7];
    piece Piece[NUM_OF_PIECES];
    short piecesInPath;
    short piecesInBase;
    short piecesInHome;
    short startingPosition;
    short homeStraightEntery;
} player;
\end{lstlisting}

\begin{description}[style=nextline, leftmargin=1.5em]
  \item[\code{color}]
    Represents the colour of the player's pieces.

  \item[\code{Piece}]
    An array holding the four pieces that belong to the player. Each element is
    an instance of the \code{piece} structure, allowing each piece's state to be
    managed individually.

  \item[\code{piecesInPath}]
    Tracks the number of pieces currently on the standard path, which helps in
    choosing which piece to move.

  \item[\code{piecesInBase}]
    Tracks how many pieces are still in the base.

  \item[\code{piecesInHome}]
    Counts how many pieces have reached home, which is essential for determining
    when a player has finished the game.

  \item[\code{startingPosition}]
    Represents each player's starting cell, so that the structure can be
    generalised without considering which player is currently playing.

  \item[\code{homeStraightEntery}]
    Represents each player's approach cell, again allowing the player structure
    to be generalised without considering which player is currently playing.
\end{description}

\subsection{Justification for the Chosen Structures}

The \code{player} structure encapsulates all attributes that define a player's
state, making it easy to manage each player's pieces collectively. The nested
\code{piece} array, in turn, allows each piece of a particular player to be
accessed and updated separately.

By maintaining counts of pieces in each state (\code{piecesInPath},
\code{piecesInBase}, \code{piecesInHome}), the structure enables quick checks
and updates without repeatedly iterating over the pieces array. In addition, by
associating each player with a starting position and a home-straight entry
cell, the same logic can be applied to every player, which keeps the code
general.

% =====================================================================
\section{Function Justification and Code Efficiency}

\begin{enumerate}[label=\arabic*.,leftmargin=*,itemsep=0.8em]

  \item \code{diceRoll()}
  \begin{itemize}[nosep]
    \item Generates a random number between 1 and 6, simulating a dice roll.
    \item \textbf{Efficiency:} $O(1)$.
  \end{itemize}

  \item \code{nextplayer(int currentPlayer)}
  \begin{itemize}[nosep]
    \item Determines the next player in sequence, ensuring the game progresses
          in cyclic order (Red $\rightarrow$ Green $\rightarrow$ Yellow
          $\rightarrow$ Blue).
    \item \textbf{Efficiency:} $O(1)$.
  \end{itemize}

  \item \code{InitializePlayer()}
  \begin{itemize}[nosep]
    \item Initialises the properties of each player, including their pieces
          and starting positions.
    \item \textbf{Efficiency:} $O(n^2)$. However, this function is called only
          once per game, so it does not affect overall efficiency.
  \end{itemize}

  \item \code{EligibleToHomeStraight(int currentPosition,}\newline\hspace*{1.5em}\code{int expectedLocation, int currentPlayer)}
  \begin{itemize}[nosep]
    \item Iterates from the current position to the expected location of a
          piece (\code{expectedLocation = currentPosition + diceRoll}) to
          determine whether the player's approach cell is passed during the
          move. If so, the piece enters the home straight instead of continuing
          along the standard path.
    \item \textbf{Efficiency:} $O(1)$, since the number of steps is bounded by
          the dice value.
  \end{itemize}

  \item \code{getPiecesInPath(int p, int pathPieces[])}
  \begin{itemize}[nosep]
    \item Returns an array containing the indexes of the pieces currently on
          the standard path. This is required to decide which piece to move
          when more than one piece is on the path. It is also used to identify
          which pieces are in the base, so that a piece can be moved to the
          starting cell when a 6 is rolled.
    \item \textbf{Efficiency:} $O(n)$.
  \end{itemize}

  \item \code{positionMaker(int index, int player, int new\_position)}
  \begin{itemize}[nosep]
    \item Sets the location of a piece that is about to enter the home
          straight from the standard path.
    \item \textbf{Efficiency:} $O(1)$.
  \end{itemize}

  \item \code{updatePosition(int p)}
  \begin{itemize}[nosep]
    \item Iterates through each piece of the given player and recalculates the
          number of pieces in the base, on the path, and at home. This ensures
          that the game state remains accurate after every move.
    \item \textbf{Efficiency:} $O(n)$.
  \end{itemize}

  \item \code{isCapture(int currentPlayer,}\newline\hspace*{1.5em}\code{int *capturedPlayer, int *capturedPiece)}
  \begin{itemize}[nosep]
    \item Takes the index of the current player so that only the pieces of the
          other three players are checked for capture. It also takes two
          pointers, to which it assigns the player index and piece index of the
          captured piece. This identifies which piece of which player must be
          returned to the base.
    \item \textbf{Efficiency:} $O(n^2)$.
  \end{itemize}

  \item \code{reset(int capturedPlayer, int capturedPiece)}
  \begin{itemize}[nosep]
    \item Called after a capture to return the captured piece to the base of
          its player.
    \item \textbf{Efficiency:} $O(1)$.
  \end{itemize}

  \item \code{movePiece(int currentPlayer, int roll)}
  \begin{itemize}[nosep]
    \item Handles all movement logic based on the dice roll, including
          capturing and entering the home straight. Encapsulating all
          movement-related logic in one function keeps the code modular and
          easy to maintain.
    \item \textbf{Efficiency:} $O(n^2)$, as it must check multiple conditions.
  \end{itemize}

  \item \code{playGame(int firstPlayer)}
  \begin{itemize}[nosep]
    \item Takes the index of the starting player, which is determined in
          \code{main}. This is the central function that drives the game,
          controlling the flow of turns and checking for a winner.
    \item \textbf{Efficiency:} The main loop runs until a winner is found, so
          its running time depends on the progress of the game.
  \end{itemize}

\end{enumerate}

% =====================================================================
\section{Summary of Time Complexities}

\begin{table}[h]
\centering
\begin{tabular}{>{\ttfamily}l l}
\toprule
\textrm{\textbf{Function}} & \textbf{Time Complexity} \\
\midrule
diceRoll               & $O(1)$ \\
nextplayer             & $O(1)$ \\
InitializePlayer       & $O(n^2)$ (called once) \\
EligibleToHomeStraight & $O(1)$ \\
getPiecesInPath        & $O(n)$ \\
positionMaker          & $O(1)$ \\
updatePosition         & $O(n)$ \\
isCapture              & $O(n^2)$ \\
reset                  & $O(1)$ \\
movePiece              & $O(n^2)$ \\
playGame               & Depends on game progress \\
\bottomrule
\end{tabular}
\caption{Time complexity of each function.}
\end{table}

\end{document}
